\section{PyTorch Model Implementation}

This section describes the PyTorch implementation of the keyword detection model, covering dataset processing pipelines, the dataset validation GUI, custom features, neural network layers, training mechanisms, and model export details.

\subsection{Dataset Curation and Restructuring Pipeline}
The training dataset is sourced from the Google Speech Commands dataset (v2). The script \texttt{Download\_Dataset.py} restructures the raw files to fit the target classes:
\begin{enumerate}
    \item \textbf{Keyword Preservation:} Directories corresponding to core keywords (\texttt{yes}, \texttt{no}, \texttt{up}, \texttt{down}) are kept.
    \item \textbf{Other Class Consolidation:} Non-target directory files (e.g., \texttt{bed}, \texttt{cat}, \texttt{happy}) are moved into a unified \texttt{dataset/other/} folder. Files are renamed using their source folder prefix to prevent naming collisions:
    \begin{equation}
        \text{dataset/bed/00f0204f\_nohash\_0.wav} \rightarrow \text{dataset/other/bed\_00f0204f\_nohash\_0.wav}
    \end{equation}
    \item \textbf{Background Noise Slicing:} Background sounds like \texttt{white\_noise.wav} and \texttt{pink\_noise.wav} are sliced into overlapping 1-second clips to act as negative samples. The number of clips is defined as:
    \begin{equation}
        N_{\text{clips}} = \lfloor \text{duration in seconds} \times 3 \rfloor
    \end{equation}
    For \texttt{doing\_the\_dishes.wav}, the slicing is overridden to generate exactly 270 clips.
\end{enumerate}

\subsection{Resolving Memory and Disk I/O Bottlenecks}
During early development, loading audio files from disk on-the-fly inside the PyTorch data loader loop led to low GPU utilization ($\sim 30\%$) and high CPU load ($70\%$--$100\%$). To resolve this disk I/O bottleneck, a RAM caching strategy was implemented in the \texttt{KeywordDataset} class.

All WAV files are read, decoded, converted to mono, and standardized to exactly $16,000$ samples (1 second at 16 kHz sample rate) during dataset initialization. The standardized tensors are cached directly in RAM:
\begin{lstlisting}[language=Python, caption={RAM Caching Dataset Implementation}, label={lst:dataset_caching}]
class KeywordDataset(Dataset):
    def __init__(self, paths, labels, config, is_training=False):
        self.paths = paths
        self.labels = labels
        self.config = config
        self.is_training = is_training
        
        self.waveforms = []
        print(f"Pre-loading {len(paths)} audio files into RAM...")
        for path in paths:
            clean_path = os.path.normpath(path)
            try:
                waveform_np, sr = sf.read(clean_path, dtype="float32")
                waveform = torch.from_numpy(waveform_np)
            except Exception as e:
                waveform = torch.zeros((1, config.NUM_SAMPLES))
            
            # Channel reduction [stereo to mono]
            if waveform.ndim == 1:
                waveform = waveform.unsqueeze(0)
            else:
                waveform = waveform.t()

            if waveform.shape[0] > 1:
                waveform = torch.mean(waveform, dim=0, keepdim=True)
                
            # Pad or crop to exactly NUM_SAMPLES
            if waveform.shape[1] > self.config.NUM_SAMPLES:
                if not self.is_training:
                    start = (waveform.shape[1] - self.config.NUM_SAMPLES) // 2
                    waveform = waveform[:, start:start + self.config.NUM_SAMPLES]
            elif waveform.shape[1] < self.config.NUM_SAMPLES:
                waveform = F.pad(waveform, (0, self.config.NUM_SAMPLES - waveform.shape[1]))
                
            self.waveforms.append(waveform)
\end{lstlisting}

During training, data augmentation (like random time shifts) is applied dynamically to the cached tensor in memory via \texttt{\_\_getitem\_\_}:
\begin{lstlisting}[language=Python, caption={In-Memory Data Augmentation}]
def __getitem__(self, idx):
    waveform = self.waveforms[idx].clone()

    if self.is_training:
        shift = random.randint(-1600, 1600)
        if shift > 0:
            waveform = F.pad(waveform, (shift, 0))[:, :-shift]
        elif shift < 0:
            waveform = F.pad(waveform, (0, -shift))[:, -shift:]

    return waveform, torch.tensor(self.labels[idx], dtype=torch.long)
\end{lstlisting}

\subsection{PyQt6 GUI Dataset Cleaning Tool}
A custom interactive desktop utility, \texttt{Wav\_File\_Cleanup.py}, was developed using PyQt6 to find silent, corrupted, or mislabeled audio files. 

\begin{enumerate}
    \item \textbf{Silent-First Sorting:} Computes the Root Mean Square (RMS) energy for all WAV files:
    \begin{equation}
        \text{RMS} = \sqrt{\frac{1}{N} \sum_{i=1}^{N} x_i^2}
    \end{equation}
    Files are sorted in ascending order of RMS, putting silent and low-energy files at the front of the review queue.
    \item \textbf{GUI Review Panel:} Plots the waveform using \texttt{pyqtgraph} and loops playback via \texttt{sounddevice}.
    \item \textbf{Action Logging:} The user review actions (\texttt{keep}, \texttt{delete}, or move to another category) are logged to a JSON file (e.g., \texttt{down\_data.json}).
    \item \textbf{Automated Cleanup:} The dataset setup utility reads these JSON logs and automatically applies the actions to the files.
\end{enumerate}

\subsection{Stratified Splits and Balancing}
The \texttt{other} category contains over 56,000 files, causing class imbalance. To resolve this:
\begin{enumerate}
    \item The \texttt{other} category is downsampled to exactly 10,000 samples.
    \item The dataset is split into training (70\%), validation (15\%), and testing (15\%) sets using stratified splits:
    \begin{lstlisting}[language=Python, caption={Stratified Splitting Implementation}]
temp_paths, test_paths, temp_labels, test_labels = train_test_split(
    initial_paths, initial_labels, test_size=0.15, stratify=initial_labels, random_state=42
)
train_paths_raw, val_paths, train_labels_raw, val_labels = train_test_split(
    temp_paths, temp_labels, test_size=0.1765, stratify=temp_labels, random_state=42
)
    \end{lstlisting}
    \item The keyword classes in the training set are oversampled to 8,000 samples each to balance them against the training samples of the \texttt{other} class.
\end{enumerate}

\subsection{Feature Extraction and Transform Math}
Rather than using raw audio or MFCCs, the final model uses Log-Mel Spectrogram features. The waveform is converted on-the-fly using:
\begin{equation}
    \text{Log-Mel}(x) = \log(\text{MelSpectrogram}(x) + 10^{-6})
\end{equation}
The transform uses a sample rate of 16 kHz, an FFT size of 400 ($25\text{ ms}$ window), a hop size of 160 ($10\text{ ms}$ window step), and 64 mel filterbank channels. The input shape $(Batch, 1, 16000)$ is mapped to $(Batch, 1, 64, 101)$.

\subsection{CNN Model Architecture}
The classification network is a 2D Convolutional Neural Network (\texttt{KeywordCNN}) with four convolutional layers.

The output spatial dimensions of a 2D convolutional layer are defined by:
\begin{equation}
    H_{\text{out}} = \left\lfloor \frac{H_{\text{in}} - K + 2P}{S} \right\rfloor + 1, \quad W_{\text{out}} = \left\lfloor \frac{W_{\text{in}} - K + 2P}{S} \right\rfloor + 1
\end{equation}
For convolutional layers with a kernel size $K=3$, padding $P=1$, and stride $S=1$:
\begin{equation}
    H_{\text{out}} = \left\lfloor \frac{H_{\text{in}} - 3 + 2(1)}{1} \right\rfloor + 1 = H_{\text{in}}
\end{equation}
Spatial downsampling is handled by Max Pooling layers with $K_{\text{pool}}=2$ and $S_{\text{pool}}=2$:
\begin{equation}
    H_{\text{pool}} = \left\lfloor \frac{H_{\text{in}}}{2} \right\rfloor
\end{equation}

\begin{table}[H]
\centering
\caption{Layer Configuration and Activation Dimensions of KeywordCNN}
\begin{tabular}{lllll}
\toprule
\textbf{Layer Name} & \textbf{Operator} & \textbf{Filters/Features} & \textbf{Parameters} & \textbf{Output Shape} \\
\midrule
Input & - & - & - & $(Batch, 1, 64, 101)$ \\
Conv 1 & 2D Conv + BN & 16 & $3\times3$, pad 1 & $(Batch, 16, 64, 101)$ \\
Pool 1 & Max Pool & - & $2\times2$, stride 2 & $(Batch, 16, 32, 50)$ \\
Conv 2 & 2D Conv + BN & 32 & $3\times3$, pad 1 & $(Batch, 32, 32, 50)$ \\
Pool 2 & Max Pool & - & $2\times2$, stride 2 & $(Batch, 32, 16, 25)$ \\
Conv 3 & 2D Conv + BN & 64 & $3\times3$, pad 1 & $(Batch, 64, 16, 25)$ \\
Pool 3 & Max Pool & - & $2\times2$, stride 2 & $(Batch, 64, 8, 12)$ \\
Conv 4 & 2D Conv + BN & 128 & $3\times3$, pad 1 & $(Batch, 128, 8, 12)$ \\
Pool 4 & Max Pool & - & $2\times2$, stride 2 & $(Batch, 128, 4, 6)$ \\
Adaptive Pool & Adaptive Avg Pool & - & Target size $4\times4$ & $(Batch, 128, 4, 4)$ \\
Flatten & Reshape & - & - & $(Batch, 2048)$ \\
FC 1 & Linear + Dropout & 128 & Dropout $p=0.3$ & $(Batch, 128)$ \\
FC 2 (Output) & Linear & $N_{\text{classes}}$ & - & $(Batch, 5)$ \\
\bottomrule
\end{tabular}
\label{tab:pytorch_cnn_dims}
\end{table}

\subsection{Training Loop, Schedulers, and Early Stopping}
The model is trained using Cross-Entropy Loss:
\begin{equation}
    \mathcal{L} = -\sum_{c=1}^{C} y_c \log(\hat{y}_c)
\end{equation}
The Adam optimizer is configured with learning rate $\eta = 0.001$. A \texttt{ReduceLROnPlateau} scheduler decreases the learning rate by a factor of 0.5 if validation accuracy does not improve for 5 epochs. Early stopping halts training if validation accuracy plateaus for 15 consecutive epochs.

\subsection{progressive Testing Iterations}
The model was optimized over 13 progressive versions:
\begin{table}[H]
\centering
\caption{Evolution of Test Accuracies Across Iterative Configurations}
\begin{tabular}{lll}
\toprule
\textbf{Script Name} & \textbf{Modification Details} & \textbf{Final Test Accuracy} \\
\midrule
\texttt{01\_specaugment.py} & Spectrogram frequency and time masking & 95.10\% -- 95.68\% \\
\texttt{02\_audio\_data\_augmentation.py} & Waveform time-shifting ($\pm 100\text{ ms}$) & 95.52\% -- 96.15\% \\
\texttt{03\_lr\_scheduler.py} & Plateau Scheduler & 95.68\% \\
\texttt{04\_batch\_normalization.py} & Batch Normalization (Bug in early version) & 53.95\% -- 64.44\% \\
\texttt{05\_weight\_decay.py} & L2 regularization ($10^{-4}$) & 94.89\% -- 95.68\% \\
\texttt{06\_increase\_filters\_and\_mfcc.py} & 4-layer CNN (128 filters), 40-coeff MFCC & 96.79\% -- 97.68\% \\
\texttt{07\_increase\_mel\_bins.py} & 128 mel bins & 89.67\% -- 94.47\% \\
\texttt{08\_increase\_dropout.py} & Dropout probability $p=0.3$ & 95.52\% \\
\texttt{09\_model\_checkpointing.py} & Best weights tracking & 95.15\% -- 95.42\% \\
\texttt{10\_combined\_best.py} & Combined enhancements & 97.95\% \\
\texttt{11\_combined\_stable.py} & Parameter standardization & 98.21\% -- 98.52\% \\
\texttt{12\_confusion\_matrix.py} & Confusion matrix evaluation & 98.58\% \\
\texttt{13\_kfold\_cross\_validation.py} & 5-Fold Cross Validation & \textbf{98.37\%} (Mean) $\pm$ \textbf{0.28\%} \\
\bottomrule
\end{tabular}
\label{tab:pytorch_testing_history}
\end{table}

\subsection{Voice Activity Detection and Confidence Filters}
In real-time environments, background noise can cause false positive detections. Two post-processing filters are applied to prevent this:

\begin{enumerate}
    \item \textbf{Voice Activity Detection (VAD):} Computes the RMS of the audio buffer. If it falls below a threshold:
    \begin{equation}
        \text{RMS} < \theta_{\text{VAD}} = 0.002
    \end{equation}
    the output is immediately overridden to the index of the \texttt{other} class.
    \item \textbf{Confidence Thresholding:} The softmax probability of the predicted class must exceed a confidence threshold to be accepted:
    \begin{equation}
        \max_c P(y=c|x) \ge \theta_{\text{conf}} = 0.85
    \end{equation}
    If the probability is below this threshold, the output is overridden to \texttt{other}.
\end{enumerate}

These filters reduced false positive detections in the live web application by approximately 50\%.

\subsection{ONNX Export}
The model is exported to ONNX format to allow real-time inference in the desktop application without running a full PyTorch environment in the frontend process. The dynamic axes are configured to handle variable batch sizes and time feature dimensions:
\begin{lstlisting}[language=Python, caption={ONNX Export Script}]
dynamic_axes = {
    'mfcc_input': {0: 'batch_size', 3: 'time'},
    'logits': {0: 'batch_size'}
}
torch.onnx.export(
    best_model,
    dummy_input,
    "PyTorch.onnx",
    export_params=True,
    opset_version=18,
    do_constant_folding=True,
    input_names=['mfcc_input'],
    output_names=['logits'],
    dynamic_axes=dynamic_axes
)
\end{lstlisting}
The exported file is saved as \texttt{PyTorch.onnx} and executed using ONNX Runtime.

\subsection{Centralized Configuration and Shared Environment}
To prevent configuration drift and coordinate operations between training notebooks, command-line dataset analysis scripts, and production servers, a centralized configuration file \texttt{config.json} is maintained in the repository root. This file defines key parameters:
\begin{lstlisting}[language=json, caption={Repository config.json Configuration Schema}]
{
  "keywords": ["yes", "no", "up", "down", "other"],
  "target_sample_rate": 16000,
  "num_samples": 16000,
  "DATA_DIR": "dataset/"
}
\end{lstlisting}

A utility library, \texttt{Utils/config\_loader.py}, provides a standardized API to query these values across Python environments:
\begin{lstlisting}[language=Python, caption={Utility Configuration Loader}]
import os
import json

def load_config():
    base_dir = os.path.dirname(os.path.abspath(__file__))
    root_dir = os.path.abspath(os.path.join(base_dir, ".."))
    config_path = os.path.join(root_dir, "config.json")
    if not os.path.exists(config_path):
        config_path = "config.json"
    with open(config_path, "r") as f:
        return json.load(f)

def get_keywords():
    return load_config()["keywords"]
\end{lstlisting}

\subsection{Dataset Specification Analysis and Verification}
Audio samples undergo pre-ingestion audit via the \texttt{Utils/analyze\_wavs.py} script. This utility uses the \texttt{soundfile} package to parse dataset files and outputs the following metrics:
\begin{itemize}
    \item \textbf{Sample Rate Audit:} Verifies that all WAV clips match the target frequency (16 kHz).
    \item \textbf{Duration Distribution:} Analyzes minimum, maximum, and average duration of clips, compiling occurrences into a histogram exported to \texttt{dataset\_distribution.png}.
    \item \textbf{Frame Distribution:} Ranks audio clip length by sample size, generating a tabular output representing clip lengths saved to \texttt{dataset\_statistics.txt}.
\end{itemize}

\subsection{Model Error Log and Failure Analysis}
To systematically debug misclassifications, the training and testing notebook logs all validation and testing mistakes. During the final evaluation pass, predicted indices that differ from the target labels are recorded to a tabular text file \texttt{Utils/PyTorch\_wronglist.txt}. This log features columns for the true category, the incorrect predicted category, and the specific WAV file:
\begin{lstlisting}[caption={Sample Failure Log Format (PyTorch\_wronglist.txt)}]
True Label   | Predicted    | File Path
------------------------------------------------------------
yes          | down         | 05b2db80_nohash_0.wav
yes          | up           | 190821dc_nohash_0.wav
yes          | other        | 3f2b358d_nohash_1.wav
other        | down         | bed_05b2db80_nohash_0.wav
\end{lstlisting}
By examining this log, developers can identify label noise or systematic errors, such as quiet keyword segments getting overridden to the \texttt{other} category by the VAD filter.

